\chapter{Suicide Risk Screening}
\section*{What Is This Test?} Suicide-risk screening asks directly about suicidal thoughts, plans, intent, access to means, and protective factors.\section*{Alternative Names} Suicide screening; self-harm risk assessment; suicide safety assessment.\section*{Principle} Brief validated questions identify risk, followed by a clinical evaluation that determines immediate safety needs.\section*{Why This Test Is Ordered} It is used when depression, self-harm, substance use, severe distress, psychosis, or warning signs are present and in selected routine settings.\section*{Primary vs Secondary Test} Screening is preliminary; a positive response requires prompt clinical assessment, not a laboratory test.\section*{How This Test Is Performed} A trained clinician asks privately about thoughts, behavior, timing, intent, means, supports, and reasons for living.\section*{What Preparation Is Needed} None. If there is immediate danger, contact emergency services or a crisis service now and do not remain alone.\section*{Understanding Results} Risk is dynamic and cannot be determined by one score; denial of thoughts does not eliminate risk when behavior or circumstances are concerning.\section*{Follow-up Tests} Safety planning, psychiatric evaluation, means restriction, substance assessment, and urgent emergency care may follow.\section*{References} \begin{enumerate}\item MedlinePlus. Suicide Risk Screening.\item Columbia Lighthouse Project. Suicide Severity Rating Scale resources.\end{enumerate}\clearpage\section*{Understanding Results and Follow-up}
The laboratory report should identify the specimen, method, units, reference interval or decision limit, and whether the result is positive, negative, abnormal, or inconclusive. Reference intervals vary with age, sex, pregnancy, specimen type, assay, and laboratory; a result outside a range is not automatically a diagnosis, and a result within a range does not always exclude disease. Discuss unexpected results with the ordering clinician, who can decide whether repeat or confirmatory testing, treatment, or referral is appropriate. Do not change medicines or delay urgent care based on an isolated result.
\section*{Regulatory Status and Authorization Date}This chapter describes a test category, not a specific commercial product. FDA, CDSCO, EU IVDR, and MHRA approval or registration dates therefore cannot be assigned without the assay manufacturer, product name, instrument, and intended use. For laboratory-developed tests, an individual agency approval date may not exist. See the Preface for the interpretation of regulatory status.
\clearpage
